\documentclass[10pt,a4paper]{amsart}
\usepackage[margin=19mm]{geometry}
\usepackage{amsmath,amssymb,mathtools}
\usepackage[scr=rsfs,cal=euler]{mathalfa}
\usepackage{xfrac}
\usepackage[unicode,hidelinks,bookmarks=false]{hyperref}
\newcommand{\ie}{i.e.\ }
\newcommand{\cf}{cf.\ }
\newcommand{\resp}{respectively\ }
\newcommand{\Subsection}{\subsection}
\providecommand{\nobreakdash}{\nobreak\hbox{-}}
\setlength{\emergencystretch}{2em}
\multlinegap0pt
\usepackage{amssymb}
\usepackage[shortlabels]{enumitem}
\newtheorem{proposition}{Proposition}
\newtheorem{theorem}{Theorem}
\newcommand{\Ax}{\textsf{Ax}}
\newcommand{\Elt}{\textsf{Elt}}
\newcommand{\Gen}{\textsf{Gen}}
\newcommand{\Line}{\textsf{Line}}
\newcommand{\MP}{\textsf{MP}}
\newcommand{\Pre}{\textsf{Pref}}
\newcommand{\Prf}{\textsf{Prf}}
\newcommand{\Prov}{\textsf{Prov}}
\newcommand{\Seq}{\textsf{Seq}}
\newcommand{\base}[1]{\mathfrak{#1}}
\newcommand{\closure}{\textsc{cl}}
\newcommand{\mbase}[1]{{\mathfrak{#1}^\ast}}
\newcommand{\pa}{\textsf{A}}
\renewcommand{\phi}{\varphi}
\newcommand{\proves}{\vdash}
\newcommand{\setClosedTerm}{\closure(\setTerm)}
\newcommand{\setTerm}{\mathcal{T}}
\newcommand{\supp}{\Vdash}
\title{On the Concept of Arithmetic Consequence}
\author{Alexander V. Gheorghiu}
\date{Source: arXiv:2603.09900v1 (CC BY 4.0)}
\begin{document}
\maketitle

\noindent\emph{Source and licence.} On the Concept of Arithmetic Consequence. Version arXiv:2603.09900v1 (CC BY 4.0), distributed by arXiv under the Creative Commons Attribution 4.0 International licence, http://creativecommons.org/licenses/by/4.0/, as recorded in the arXiv licence metadata for this exact version and shown on the abstract page https://arxiv.org/abs/2603.09900v1. Full licence notice: \texttt{licenses/arxiv-2603.09900-CC-BY-4.0.txt}.\par\smallskip

\noindent\textit{Editorial scope.} Counted unit: the article's theorem thm:reflection (source0.tex lines 765--767). The article's complete original proof of it is delivered with it (source0.tex lines 768--878, one \texttt{proof} environment of 111 lines). No mathematics is written, repaired or re-derived: the statement and the proof are the article's own lines, in the article's own notation, and no retained line is substituted or re-flowed.  Retained uncounted as the source-local context the proof needs: lines 356--358; lines 662--671; lines 717--721; lines 742--752.  Nothing was omitted from any retained span, so the package carries no omission record.  No retained line contains a citation, so the wrapper carries no bibliography and no bibliography entry.  No retained line contains a figure environment, an image-inclusion command or drawing code, so no image is delivered and the package declares no asset dependency.  Editorial formatting only, in addition: an AMS \texttt{amsart} preamble that restates the article's own declarations and macros used by the retained lines, namely the environments \texttt{proposition}, \texttt{theorem} on separate counters, together with the standard packages those lines need.  The retained environments print in this document as Proposition 1, Theorem 1 in order of appearance, whereas the article prints the same items with the numbering of its own full document.  The complete native source of the article as distributed in the arXiv e-print for this exact version is bundled unchanged as \texttt{sources/196097/source0.tex}.
\begin{proposition} \label{prop:local-soundness}
Let $\supp_{\base{B}} \gamma$ for $\gamma \in \Gamma$. If $\vdash_{\Gamma} \phi$, then $\supp_{\base{B}} \phi$
\end{proposition}

\begin{proposition}\label{prop:maxiconsistent-disjunction}
If $\mbase{B}$ is maxiconsistent, then:
\begin{itemize}
    \item $\supp_{\mbase{B}} \neg \varphi$ iff $\not\supp_{\mbase{B}} \varphi$;
    \item $\supp_{\mbase{B}} (\varphi \lor \psi)$ iff $\supp_{\mbase{B}} \varphi$
        or $\supp_{\mbase{B}} \psi$;
    \item $\supp_{\mbase{B}} \exists x\,\varphi$ iff $\supp_{\mbase{B}}
        \varphi[x \mapsto t]$ for some $t \in \setClosedTerm$.
\end{itemize}
\end{proposition}

\begin{proposition}\label{prop:maxiconsistent-extension}
If $\not\supp_{\base{B}} \varphi$, then 
there exists a maxiconsistent base $\mbase{B} \supseteq \base{B}$ 
such that $\not\supp_{\mbase{B}} \varphi$.
\end{proposition}

\begin{proposition}\label{prop:equationally-complete}
    Let $\pa$ be sufficiently strong $\base{B}$ be a consistent base that supports all
axioms of $\pa$. For any $\Delta_0$-formula $\delta$, if
\[
    \supp_{\base{B}} \delta,
\]
then
\[
    \proves_{\pa} \delta.
\] 
\end{proposition}

\begin{theorem}\label{thm:reflection}
    For any $\varphi$, $\pa \supp \Prov_{\pa}(\ulcorner \varphi \urcorner) \to \varphi$.
\end{theorem}

\begin{proof}
    The statement is equivalent to $\pa, \Prov_{\pa}(\ulcorner \varphi \urcorner) \supp \varphi$. Assume, for contradiction, that this is not the case: $\pa, \Prov_{\pa}(\ulcorner \varphi \urcorner) \not \supp \varphi$. Then there is a base $\base{B}$ such that:
    \begin{itemize}
        \item $\supp_{\base{B}} \alpha \in \pa$,
        \item $\supp_{\base{B}}  \Prov_{\pa}(\ulcorner \varphi \urcorner)$, but
        \item $\not \supp_{\base{B}} \varphi$.
    \end{itemize}
    Without loss of generality, by Proposition~\ref{prop:maxiconsistent-extension}, $\base{B}$ is maxiconsistent. We show that from the first two we get $\supp_{\base{B}} \varphi$, contradicting the third. \medskip

By Proposition~\ref{prop:maxiconsistent-disjunction}, there is a term $p$ such that 
\[
\supp_{\base{B}}  \Prf_{\pa}(p,\ulcorner \varphi \urcorner)
\]
Unfolding the definition of $\Prf$,
this abbreviates
\[
\supp_{\base{B}}\exists n\Big(
\Seq(p,n)\wedge n>0 \wedge \Elt(p,n-1,\ulcorner\varphi\urcorner)\wedge \forall i<n\,\Line(p,i)
\Big).
\]
By Proposition~\ref{prop:maxiconsistent-disjunction} again and the clause for $\land$, there is some $n$ such that
$\base{B}$ supports each conjunct:
\begin{enumerate}
\item $\supp_{\base{B}}\Seq(p,\bar n)$,
\item $\supp_{\base{B}}\bar n>0$,
\item $\supp_{\base{B}}\Elt(p,\overline{n-1},\ulcorner\varphi\urcorner)$,
\item $\supp_{\base{B}}\forall i<\bar n\,\Line(p,i)$.
\end{enumerate}

We show $\supp_{\base{B}}\varphi$ by induction on $n$ --- that is, the length of the coded proof $p$.

\smallskip
\noindent\textbf{Base case:} $n=1$. From (3) we obtain that the $0$th line of $\varphi$ is $\ulcorner \varphi \urcorner$,
\[
\supp_{\base{B}} \Elt(p,0,\ulcorner\varphi\urcorner) \qquad \text{and} \qquad \supp_{\base{B}} \Line(p,0)
\]
From Proposition~\ref{prop:maxiconsistent-disjunction} on $\Line(p,0)$ we get that 
\[
\Elt(p,0,t) \qquad \text{ and } \supp_{\base{B}} \Ax(t)
\]
for some term $t$. As $\Ax$ is $\Delta_0$, from Proposition~\ref{prop:equationally-complete}, we have $\vdash_{\pa} \Ax(t)$. 

As $\Elt(p,0,\ulcorner\varphi\urcorner)$ is also $\Delta_0$, from Proposition~\ref{prop:equationally-complete} we also get $\vdash_{\pa} \Elt(p,0,\ulcorner\varphi\urcorner)$. By functionality of sequence encoding (\textbf{Seq}), we have that $t = \ulcorner\varphi\urcorner$. By correctness of the rules (\textbf{Rule}), we obtain $\varphi = \alpha$ for some $\alpha \in \pa$. Hence, as $\base{B}$ supports the axioms of $\pa$, we have $\supp_{\base{B}} \varphi$.

\smallskip
\noindent\textbf{Inductive step:} Assume $n>1$ and, as \emph{induction hypothesis} (IH), that the claim holds for all
proofs of length $k<n$; that is, if
\[
\supp_{\base{B}} \Prf_{\pa}(p', \ulcorner \varphi' \urcorner) \qquad \text{ and } \qquad \supp_{\base{B}} \Seq(p',\bar{k}) \text{ with } k<n, 
\]
then 
\[
\supp_{\base{B}} \psi
\]
From (4) we obtain
$\supp_{\base{B}}\Line(p,\overline{n-1})$. 
Unfolding the definition and using Proposition~\ref{prop:maxiconsistent-disjunction}, there are terms
$i,j$ and $y_1,y_2,y$ such that
\[
\supp_{\base{B}}\Elt(p,\overline{n-1},y)
\]
and one of the following holds:
\begin{itemize}
\item $\supp_{\base{B}}\Ax(y)$; or
\item $\supp_{\base{B}}
\Elt(p,i,y_1)\wedge \Elt(p,j,y_2)\wedge \MP(y_1,y_2,y)$; or
\item $\supp_{\base{B}}
\Elt(p,j,y_1)\wedge \Gen(y_1,y)$.
\end{itemize}
As these are all $\Delta_0$-formulae, by Proposition~\ref{prop:equationally-complete}, we have that
\[
\vdash_{\pa} \Elt(p,\overline{n-1},y)
\]
and one of the following holds:
\begin{itemize}
\item $\vdash_{\pa}\Ax(y)$; or
\item $\vdash_{\pa}\Elt(p,i,y_1)\wedge \Elt(p,j,y_2)\wedge \MP(y_1,y_2,y)$; or
\item $\vdash_{\pa}
\Elt(p,j,y_1)\wedge \Gen(y_1,y)$.
\end{itemize}
From $\vdash_{\pa} \Elt(p,\overline{n-1},y)$ and (3) it follows from (\textbf{Seq}) that $y= \ulcorner \varphi \urcorner$. We now consider each disjunct in turn and show that $\supp_{\base{B}} \phi$ holds: 
\begin{itemize}
    \item We reason exactly as in the base case.
    \item By invertibility of conjunction, have that
\[
\vdash_{\pa} \Elt(p, i, y_1),\qquad \vdash_{\pa}\Elt(p, j, y_2),\qquad \vdash_{\pa}\MP( y_1, y_2,\ulcorner\varphi\urcorner).
\]
By the correctness of the encoding (\textbf{Rule}), there is a formula $\psi$ such that
$ y_1=\ulcorner\psi\urcorner$ and $ y_2=\ulcorner\psi\to\varphi\urcorner$. Let
$p_1:=\Pre(p, i+\bar{1})$ and $p_2:=\Pre(p, j+\bar{1})$; by the prefix closure on proofs (\textbf{Pref}), 
\[
\vdash_{\pa} \Prf(p_1,\ulcorner\psi\urcorner)
\quad\text{and}\quad
\vdash_{\pa} \Prf(p_2,\ulcorner\psi\to\varphi\urcorner)
\]
By local soundness (Proposition~\ref{prop:local-soundness}),
\[
\supp_{\base{B}}  \Prf(p_1,\ulcorner\psi\urcorner)
\quad\text{and}\quad
\supp_{\base{B}} \Prf(p_2,\ulcorner\psi\to\varphi\urcorner)
\]
By the length of prefixes ($\textbf{Len}$), the length of $p_1$ and $p_2$ are less than the length of $p$. Applying IH to each, we obtain
$\supp_{\base{B}}  \psi$ and $\supp_{\base{B}} \psi\to\varphi$. The desired result obtains by \emph{modus ponens}.

\item This is analogous to the previous case. Using \textbf{Rule}, there is some $\psi$ with
$y_1=\ulcorner\psi\urcorner$ such that $\varphi=\forall x \psi$.  From \textbf{Pref} and \textbf{Len}, we obtain a 
proof of $\psi$. Using Proposition~\ref{prop:local-soundness}, we move back to $\base{B}$. The induction hypothesis yields $\supp_{\base{B}}\psi$; and the
generalization yields $\supp_{\base{B}}\varphi$.
\end{itemize}
This completes the induction.
\end{proof}

\end{document}
